\documentclass[11pt]{article}

\usepackage{preprint}

%% Math packages
\usepackage{amsmath, amsthm, amssymb, amsfonts}
\usepackage{bm}
\usepackage{mathtools}

\usepackage[style=authoryear-comp,maxcitenames=2,uniquename=false,giveninits=true,uniquelist=false,maxbibnames=99,sortcites=true,sorting=nyt,doi=true,isbn=false,url=false,hyperref=true,date=year]{biblatex}
\renewbibmacro{in:}{%
  \ifboolexpr{test {\ifentrytype{article}} or test {\ifentrytype{inproceedings}}}{}%
  {\printtext{\bibstring{in}\intitlepunct}}%
}
\AtEveryBibitem{\clearlist{language}\clearfield{note}\clearfield{abstract}}
\addbibresource{linear_unmixing_note.bib}

\usepackage[utf8]{inputenc}
\usepackage[T1]{fontenc}
\usepackage{xcolor}
\usepackage[colorlinks=true, linkcolor=magenta, urlcolor=blue, citecolor=cyan, anchorcolor=black]{hyperref}
\usepackage{booktabs}
\usepackage{microtype}
\usepackage{graphicx}
\usepackage{xspace}
\usepackage{titlesec}
\titlespacing\section{0pt}{12pt plus 3pt minus 3pt}{1pt plus 1pt minus 1pt}
\titlespacing\subsection{0pt}{10pt plus 3pt minus 3pt}{1pt plus 1pt minus 1pt}

\usepackage{tikz}
\usetikzlibrary{arrows.meta}

%% preprint.sty hardcodes the parent paper's title in the running header.
\lhead{\scshape Note -- Linear least-squares unmixing on a sample network}

%% Notation, following the parent preprint
\newcommand{\cvec}{\ensuremath{\mathbf{c}}\xspace}
\newcommand{\dvec}{\ensuremath{\mathbf{d}}\xspace}
\newcommand{\qvec}{\ensuremath{\mathbf{q}}\xspace}
\newcommand{\Mmat}{\ensuremath{\mathbf{M}}\xspace}
\newcommand{\Rmat}{\ensuremath{\mathbf{R}}\xspace}
\newcommand{\Pmat}{\ensuremath{\mathbf{P}}\xspace}
\newcommand{\Amat}{\ensuremath{\mathbf{A}}\xspace}
\newcommand{\Wmat}{\ensuremath{\mathbf{W}}\xspace}
\newcommand{\Imat}{\ensuremath{\mathbf{I}}\xspace}
\newcommand{\onevec}{\ensuremath{\mathbf{1}}\xspace}
\newcommand{\Cd}{\ensuremath{\mathbf{C}_{\mathbf{d}}}\xspace}
\newcommand{\Cc}{\ensuremath{\mathbf{C}_{\mathbf{c}}}\xspace}
\newcommand{\Cdhat}{\ensuremath{\mathbf{C}_{\hat{\mathbf{d}}}}\xspace}
\newcommand{\Reals}{\ensuremath{\mathbb{R}}\xspace}
\newcommand{\Expect}{\ensuremath{\mathbb{E}}\xspace}
\newcommand{\transpose}{\ensuremath{^{\mathsf{T}}}}

\DeclareMathOperator*{\argmin}{arg\,min}

\newtheorem{theorem}{Proposition}
\newtheorem{lemma}[theorem]{Lemma}
\newtheorem{corollary}[theorem]{Corollary}
\theoremstyle{definition}
\newtheorem{remark}[theorem]{Remark}

\title{Linear least-squares unmixing on a sample network:\\derivation, regularization, and error propagation}

\usepackage{authblk}
\renewcommand*{\Authfont}{\bfseries}
\author[1]{Alex G. Lipp}
\affil[1]{Department of Earth Sciences, University College London}
\date{}

\begin{document}
\maketitle

\begin{abstract}
\noindent
\textcite{barnes_using_2024} recover source concentrations on a river sample network by
convex optimization, penalising \textit{relative} differences between predicted and observed
concentrations. Their Appendix~A observes that penalising \textit{absolute} differences
instead reduces the problem to a linear least-squares matrix inversion. This note works that
observation out in full. We show that the mixing matrix $\Mmat$ is square, row-stochastic and
triangular, hence always invertible, and we give its inverse in closed form as a local
differencing stencil on the network. We then extend the formulation in three directions that
the appendix does not cover: first-order tracer decay, weighting by the data covariance
(generalised least squares), and Tikhonov regularization of the model variance. We prove that
the regularized problem has a unique solution for every regularization strength, derive the
resulting linear estimator, and obtain closed-form covariances for both the recovered sources
and the modelled observations. We distinguish carefully between the propagated covariance and
the Bayesian posterior covariance, which coincide only when the problem is unregularized, and
give the exact identity relating them. Everything is derived from first principles so that it
can be checked line by line.
\end{abstract}

\vspace{3mm}

\section{Scope and notation}

We adopt the notation of \textcite{barnes_using_2024} throughout. The sample network is a
directed acyclic graph $G = (N, E)$ whose $n$ nodes are sample sites and whose edges point
\textit{downstream}. Each node $i$ carries:

\begin{itemize}
  \item $a_i$ [L$^2$], the area of the sub-basin uniquely defined by sample site $i$;
  \item $\phi_i$ [M\,L$^{-2}$\,T$^{-1}$], the material export rate of that sub-basin;
  \item $q_i = a_i \phi_i > 0$ [M\,T$^{-1}$], the material flux the sub-basin generates;
  \item $c_i$, the source concentration of the tracer in material exported by sub-basin $i$
        --- the unknown;
  \item $d_i$, the concentration observed at the sample site --- the datum.
\end{itemize}

\noindent Write $U_i \subseteq N$ for the set of nodes upstream of $i$ \textit{including $i$
itself}, and $A(i)$ for the set of immediate upstream neighbours of $i$ (its predecessors in
$G$). We make one structural assumption throughout, which the D8 flow model guarantees:

\begin{quote}
\textbf{(A1)} Every node has at most one downstream neighbour, so $G$ is a forest of
in-trees.
\end{quote}

\noindent Assumption (A1) is what makes the sets $\{U_p\}_{p \in A(i)}$ pairwise disjoint,
which we use in Proposition~\ref{prop:inverse}. It does not require the network to be
connected: disjoint river basins may be solved together, exactly as in the parent study.

Two derived quantities recur. The \textbf{total flux} passing site $i$ is
\begin{equation}
  Q_i \;=\; \sum_{j \in U_i} q_j ,
  \label{eq:Q}
\end{equation}
and, for $j \in U_i$, the \textbf{path attenuation} from $j$ to $i$ is
\begin{equation}
  \alpha_{j \to i} \;=\; \prod_{e \in \mathrm{path}(j \to i)} \exp\!\left( -k_{e} L_{e} \right),
  \label{eq:alpha}
\end{equation}
the product over edges of the unique path from $j$ to $i$, where $L_e$ is the flow-path length
of edge $e$ and $k_e$ the first-order decay constant applied along it. By (A1) that path is
unique, so \eqref{eq:alpha} is well defined, and $\alpha_{i \to i} = 1$ (empty product). A
\textbf{conservative} tracer has $k_e = 0$ for all $e$, hence $\alpha_{j \to i} = 1$
throughout; this is the case treated in Appendix~A of the parent paper.

\section{The forward model as a matrix}

\subsection{Assembling \texorpdfstring{$\Mmat$}{M}}

Conservation of tracer at steady state gives the observation at site $i$ as a flux-weighted
mixture of every source upstream of it, attenuated by decay along the way:
\begin{equation}
  d_i \;=\; \frac{\displaystyle\sum_{j \in U_i} \alpha_{j \to i}\, q_j\, c_j}
                 {\displaystyle\sum_{j \in U_i} q_j}
        \;=\; \frac{1}{Q_i} \sum_{j \in U_i} \alpha_{j \to i}\, q_j\, c_j .
  \label{eq:forward}
\end{equation}

\noindent Note that decay enters the numerator only. Tracer is lost in transit, but the
\textit{carrier} --- the water or sediment --- is not, so the denominator is the undecayed
total flux $Q_i$. Setting every $\alpha = 1$ recovers Equation~3 of \textcite{barnes_using_2024}.

Following the parent appendix, define the $0$--$1$ matrix $\bm{\Theta}$ by $\Theta_{ij} = 1$
if and only if $j \in U_i$; that is, $\bm{\Theta}$ is the path matrix of $G$ with a self-loop
added at every node. Equation~\eqref{eq:forward} is then linear in \cvec, with
\begin{equation}
  \boxed{\;
  M_{ij} \;=\; \frac{\Theta_{ij}\, \alpha_{j \to i}\, q_j}
                    {\sum_{k=1}^{n} \Theta_{ik}\, q_k}
             \;=\; \frac{\Theta_{ij}\, \alpha_{j \to i}\, q_j}{Q_i} ,
  \qquad \dvec = \Mmat \cvec .
  \;}
  \label{eq:M}
\end{equation}

\noindent For a conservative tracer this is exactly Equation~11 of the parent appendix.

\subsection{Structure of \texorpdfstring{$\Mmat$}{M}}

$\Mmat$ is $n \times n$: the sample network defines exactly one sub-basin per sample site, so
there are as many unknowns as data. This squareness is the pivot on which everything below
turns, and it is worth stating plainly that it is a modelling choice, not a coincidence --- it
is what it means to attribute one source concentration to each sampled sub-basin.

\begin{lemma}[Row sums]\label{lem:stochastic}
Every row of $\Mmat$ is non-negative with $\sum_j M_{ij} \le 1$, with equality for all $i$ if
and only if the tracer is conservative.
\end{lemma}

\begin{proof}
Non-negativity is immediate from \eqref{eq:M} since $q_j > 0$, $\alpha > 0$, $\Theta \in
\{0,1\}$. For the sum,
\[
  \sum_{j=1}^n M_{ij}
  = \frac{1}{Q_i}\sum_{j \in U_i} \alpha_{j\to i}\, q_j
  \;\le\; \frac{1}{Q_i} \sum_{j \in U_i} q_j
  \;=\; \frac{Q_i}{Q_i} = 1 ,
\]
using $\alpha_{j \to i} \le 1$, which holds because every $k_e L_e \ge 0$. Equality for a
given $i$ requires $\alpha_{j \to i} = 1$ for all $j \in U_i$, i.e.\ no decay anywhere
upstream of $i$; equality for all $i$ requires it everywhere.
\end{proof}

\noindent For a conservative tracer, then, $\Mmat$ is \textbf{row-stochastic}: each observation
is a convex combination of the sources upstream of it. This is the formal statement of ``a
river sample is a mixture of its upstream sources'', and it has a useful consequence,
$\Mmat \onevec = \onevec$, which we use in Section~\ref{sec:limits}.

\begin{lemma}[Triangularity]\label{lem:triangular}
Let $\sigma$ be any topological ordering of $G$, so that $\sigma(u) < \sigma(v)$ whenever
$(u,v) \in E$. Permuting rows and columns of $\Mmat$ by $\sigma$ makes it lower triangular,
with diagonal entries
\begin{equation}
  M_{ii} \;=\; \frac{q_i}{Q_i} \;\in\; (0, 1] .
  \label{eq:diagonal}
\end{equation}
\end{lemma}

\begin{proof}
$M_{ij} \ne 0$ requires $\Theta_{ij} = 1$, i.e.\ $j \in U_i$. If $j \ne i$ then a directed
path runs from $j$ to $i$, so $\sigma(j) < \sigma(i)$. Hence $M_{ij} \ne 0 \Rightarrow
\sigma(j) \le \sigma(i)$, which is lower triangularity in the permuted ordering. The diagonal
follows from \eqref{eq:M} with $j = i$, using $\Theta_{ii} = 1$ and $\alpha_{i \to i} = 1$; it
is positive because $q_i > 0$, and at most $1$ because $q_i \le Q_i$, with $M_{ii} = 1$ exactly
when $i$ is a leaf (headwater) node.
\end{proof}

\begin{theorem}[Invertibility]\label{prop:invertible}
$\Mmat$ is invertible, with
\begin{equation}
  \left| \det \Mmat \right| \;=\; \prod_{i=1}^{n} \frac{q_i}{Q_i} \;>\; 0 .
\end{equation}
\end{theorem}

\begin{proof}
Let $\bm{\Pi}$ be the permutation matrix of the topological order $\sigma$. By
Lemma~\ref{lem:triangular}, $\bm{\Pi}\Mmat\bm{\Pi}\transpose$ is lower triangular, so its
determinant is the product of its diagonal entries, $\prod_i q_i / Q_i$, which is strictly
positive because every $q_i > 0$ and every $Q_i < \infty$. Since $|\det \bm{\Pi}| = 1$,
$|\det \Mmat| = \prod_i q_i/Q_i > 0$, and a matrix with non-zero determinant is invertible.
\end{proof}

\begin{remark}
Invertibility is unaffected by decay: $\alpha$ appears only off the diagonal, so
\eqref{eq:diagonal} is untouched. It is also unaffected by the export rates $\phi$, provided
they are strictly positive. It \textit{would} fail if some sub-basin had zero area or zero
export rate, i.e.\ if two sample sites were co-located --- then $q_i = 0$, the corresponding
source is not identifiable, and the model is degenerate. This is worth checking in practice.
\end{remark}

\section{The inverse in closed form}
\label{sec:inverse}

Proposition~\ref{prop:invertible} says $\cvec = \Mmat^{-1}\dvec$ exists and is unique. We now show
it need never be computed by matrix inversion: it is a local differencing stencil on the
network, computable in one sweep.

It is convenient to work with the \textbf{tracer flux} passing site $i$,
\begin{equation}
  T_i \;=\; Q_i d_i \;=\; \sum_{j \in U_i} \alpha_{j \to i}\, q_j\, c_j ,
  \label{eq:T}
\end{equation}
which is just \eqref{eq:forward} cleared of its denominator. $T_i$ has units of tracer mass per
time: it is the amount of tracer passing the sample site per unit time.

\begin{theorem}[Closed-form inverse]\label{prop:inverse}
For every node $i$,
\begin{equation}
  \boxed{\;
  c_i \;=\; \frac{1}{q_i}\left( Q_i d_i \;-\; \sum_{p \in A(i)} \alpha_{p \to i}\, Q_p d_p \right) .
  \;}
  \label{eq:closedform}
\end{equation}
\end{theorem}

\begin{proof}
Every node in $U_i$ other than $i$ lies upstream of exactly one immediate neighbour $p \in
A(i)$, and by (A1) the sets $U_p$ for distinct $p \in A(i)$ are pairwise disjoint. Hence
\[
  U_i \;=\; \{i\} \;\sqcup\; \bigsqcup_{p \in A(i)} U_p ,
\]
a disjoint union. Now fix $p \in A(i)$ and $j \in U_p$. The unique path $j \to i$ is the path
$j \to p$ followed by the single edge $p \to i$, so the attenuation factorises:
\begin{equation}
  \alpha_{j \to i} \;=\; \alpha_{j \to p}\;\alpha_{p \to i} .
  \label{eq:factorise}
\end{equation}
Substituting the partition and \eqref{eq:factorise} into \eqref{eq:T}, and using
$\alpha_{i \to i} = 1$ for the $j = i$ term,
\begin{align}
  T_i
  &= q_i c_i \;+\; \sum_{p \in A(i)} \; \sum_{j \in U_p} \alpha_{j \to i}\, q_j c_j \nonumber\\
  &= q_i c_i \;+\; \sum_{p \in A(i)} \alpha_{p \to i} \underbrace{\sum_{j \in U_p} \alpha_{j \to p}\, q_j c_j}_{= \,T_p} \nonumber\\
  &= q_i c_i \;+\; \sum_{p \in A(i)} \alpha_{p \to i}\, T_p .
  \label{eq:balance}
\end{align}
Equation~\eqref{eq:balance} is a statement of tracer conservation at node $i$: the tracer
leaving equals the tracer generated locally plus the (attenuated) tracer delivered by the
immediate tributaries. Solving for $c_i$ and substituting $T = Qd$ gives
\eqref{eq:closedform}. Since $q_i > 0$ the division is legitimate.
\end{proof}

\begin{corollary}[Sparsity of $\Mmat^{-1}$]
\begin{equation}
  \left(\Mmat^{-1}\right)_{ij} =
  \begin{dcases}
    \dfrac{Q_i}{q_i}, & j = i, \\[6pt]
    -\dfrac{\alpha_{j \to i}\, Q_j}{q_i}, & j \in A(i), \\[6pt]
    0, & \text{otherwise.}
  \end{dcases}
  \label{eq:Minv}
\end{equation}
$\Mmat^{-1}$ therefore has at most $n + |E| \le 2n - 1$ non-zero entries, against up to
$n(n+1)/2$ in $\Mmat$ itself. Equation~\eqref{eq:closedform} evaluates in $O(n + |E|)$ time.
\end{corollary}

\noindent It is worth pausing on how much structure this reveals. $\Mmat$ is dense and global
--- every downstream observation depends on every source above it --- yet its inverse is
sparse and purely local. Recovering a source concentration requires only the observation at
that site and the observations at its immediate tributaries. All the long-range mixing cancels.

\begin{corollary}[Worked example]
For the six-node network of \textcite[Figure~2]{barnes_using_2024}, with $A(4) = \{1,2\}$,
$A(5) = \{3\}$, $A(6) = \{4,5\}$ and no decay, Equation~\eqref{eq:closedform} gives
\begin{align*}
  c_1 &= d_1, \qquad c_2 = d_2, \qquad c_3 = d_3, \\
  c_4 &= \frac{Q_4 d_4 - Q_1 d_1 - Q_2 d_2}{q_4}, \\
  c_5 &= \frac{Q_5 d_5 - Q_3 d_3}{q_5}, \\
  c_6 &= \frac{Q_6 d_6 - Q_4 d_4 - Q_5 d_5}{q_6},
\end{align*}
with $Q_1 = q_1$, $Q_2 = q_2$, $Q_3 = q_3$, $Q_4 = q_1 + q_2 + q_4$, $Q_5 = q_3 + q_5$ and
$Q_6 = \sum_{i=1}^{6} q_i$. Headwater sites are recovered exactly and trivially; interior sites
are recovered by differencing.
\end{corollary}

\subsection{Noise amplification}
\label{sec:amplification}

Differencing is the source of all the difficulty. Suppose the observations carry independent
errors of standard deviation $\sigma_i$. Propagating them through \eqref{eq:closedform}
(anticipating Section~\ref{sec:propagation}, and using independence),
\begin{equation}
  \operatorname{Var}(c_i)
  = \frac{Q_i^2 \sigma_i^2 + \sum_{p \in A(i)} \alpha_{p\to i}^2 Q_p^2 \sigma_p^2}{q_i^2} .
  \label{eq:varc}
\end{equation}
Every term is scaled by $q_i^{-2}$. Define the \textbf{amplification factor}
\begin{equation}
  \kappa_i \;=\; \frac{Q_i}{q_i} \;=\; \frac{1}{M_{ii}} \;\ge\; 1 ,
  \label{eq:kappa}
\end{equation}
the ratio of the total flux passing site $i$ to the flux its own sub-basin generates. If
site $i$ sits immediately downstream of another site, with little incremental drainage area
between them, then $q_i \ll Q_i$, $\kappa_i$ is large, and \eqref{eq:closedform} subtracts two
nearly equal large numbers to obtain a small one. This is catastrophic cancellation, and it is
an intrinsic property of the sampling design, not of the algorithm: the data simply do not
constrain that sub-basin well. In practice $\kappa_i$, being free to compute, is the single
most useful diagnostic of whether a linear inversion is worth attempting on a given network.

\begin{corollary}[When positivity binds]\label{cor:positivity}
Since concentrations must be non-negative, the unregularized solution is physically admissible
at node $i$ if and only if
\begin{equation}
  Q_i d_i \;\ge\; \sum_{p \in A(i)} \alpha_{p \to i} \, Q_p d_p ,
  \label{eq:positivity}
\end{equation}
that is, if and only if at least as much tracer flux leaves site $i$ as its tributaries deliver
to it. A violation is a direct falsification of the forward model \eqref{eq:forward} by the
data, whether through measurement error, non-conservative behaviour, or mis-specified fluxes.
\end{corollary}

\section{Why the unregularized problem is degenerate}
\label{sec:degenerate}

Appendix~A of \textcite{barnes_using_2024} poses the recovery as a constrained least-squares
problem,
\begin{equation}
  \underset{\cvec \in \Reals^n}{\text{minimize}} \;\; \lVert \Mmat\cvec - \dvec \rVert_2
  \qquad \text{subject to} \quad c_i \ge 0 .
  \label{eq:appendixproblem}
\end{equation}
Proposition~\ref{prop:invertible} shows that, absent the constraint, this is not really a
least-squares problem at all.

\begin{theorem}[Degeneracy]\label{prop:degenerate}
Let $\Wmat$ be any symmetric positive definite weight matrix. Then
\begin{equation}
  \argmin_{\cvec \in \Reals^n} \; (\Mmat\cvec - \dvec)\transpose \Wmat (\Mmat\cvec - \dvec)
  \;=\; \Mmat^{-1}\dvec ,
\end{equation}
independently of $\Wmat$, and the residual at the optimum is exactly zero.
\end{theorem}

\begin{proof}
The objective is non-negative for every \cvec, and $\cvec = \Mmat^{-1}\dvec$ --- which exists
by Proposition~\ref{prop:invertible} --- attains the value $0$. Hence it is a global minimiser, and
it is the unique one because $\Wmat \succ 0$ forces the residual to vanish, and $\Mmat$ is
injective. Algebraically, the generalised least-squares estimator collapses:
\[
  (\Mmat\transpose\Wmat\Mmat)^{-1}\Mmat\transpose\Wmat
  = \Mmat^{-1}\Wmat^{-1}\Mmat^{-\mathsf{T}}\Mmat\transpose\Wmat
  = \Mmat^{-1} . \qedhere
\]
\end{proof}

\noindent Three consequences follow, and they set up the rest of this note.

\begin{enumerate}
  \item \textbf{Weighting is inert when $\lambda = 0$.} No choice of $\Wmat$ --- including
        $\Wmat = \Cd^{-1}$, the statistically correct one --- changes the answer. There is
        nothing to trade off when the model can fit every datum exactly. Weighting only
        becomes consequential once a competing term is added to the objective
        (Section~\ref{sec:regularization}).
  \item \textbf{The problem interpolates the noise.} A zero residual means the estimate
        reproduces the observations exactly, errors included. Combined with the amplification
        of Section~\ref{sec:amplification}, this is why the raw inversion is unusable on
        realistic data.
  \item \textbf{The constraint is the only thing that can bite.} By
        Corollary~\ref{cor:positivity}, \eqref{eq:appendixproblem} departs from
        $\Mmat^{-1}\dvec$ precisely when \eqref{eq:positivity} fails somewhere.
\end{enumerate}

\begin{remark}[Bounds]
The parent appendix writes the constraint as $0 < c_i < 1$, appropriate for a tracer measured
as a mass fraction. Two remarks. First, a strict inequality cannot be imposed in a convex
program --- the feasible set would not be closed and the infimum need not be attained --- so
one uses $c_i \ge 0$. Second, the upper bound is unit-dependent ($1$ for a fraction, $10^6$ for
mg\,kg$^{-1}$) whereas the lower bound is not. Since $\Mmat$ is row-substochastic, an estimate
exceeding the physical ceiling indicates that the model has already failed, and is more useful
surfaced as a diagnostic than silently clipped. Below we impose $c_i \ge 0$ only.
\end{remark}

\section{Scale invariance and normalisation}
\label{sec:scale}

\begin{theorem}[Homogeneity]\label{prop:homogeneity}
The unregularized solution map $\dvec \mapsto \cvec$ is homogeneous of degree one:
$\cvec(t\dvec) = t\,\cvec(\dvec)$ for any $t > 0$. Moreover $\Mmat$ is invariant under a common
rescaling of the fluxes, $q_j \mapsto sq_j$.
\end{theorem}

\begin{proof}
The first claim is immediate from linearity, $\cvec = \Mmat^{-1}\dvec$. For the second,
replacing $q_j$ by $sq_j$ in \eqref{eq:M} multiplies numerator and denominator by $s$.
\end{proof}

\noindent Homogeneity licenses the mean normalisation used in the implementation. Writing
$\bar{d} = n^{-1}\sum_i d_i$ and solving with $\tilde{\dvec} = \dvec / \bar{d}$, the recovered
sources are $\cvec = \bar{d}\,\tilde{\cvec}$ exactly. This changes nothing mathematically; its
purpose is numerical, keeping all quantities near unity so that solver tolerances --- which are
partly absolute --- behave uniformly across tracers whose concentrations may differ by many
orders of magnitude. It has the further benefit of rendering the regularization strength
introduced below dimensionless, and hence transferable between tracers and datasets.

An identical argument applies to the sub-basin areas, which the implementation divides by their
mean, and to covariances, which by Proposition~\ref{prop:homogeneity} and
Section~\ref{sec:propagation} scale as $\bar{d}^{\,2}$ and so may equivalently be computed in
the original units.

\section{Weighting by the data covariance}
\label{sec:weighting}

Let $\Cd \succ 0$ be the covariance of the observation errors. The statistically natural misfit
is the Mahalanobis norm of the residual, giving generalised least squares
\parencite{menke_geophysical_2012, aster_parameter_2018}:
\begin{equation}
  \Phi(\cvec) \;=\; (\Mmat\cvec - \dvec)\transpose \Cd^{-1} (\Mmat\cvec - \dvec) .
  \label{eq:gls}
\end{equation}
Numerically one never forms $\Cd^{-1}$. Take the Cholesky factorisation $\Cd = \mathbf{L}
\mathbf{L}\transpose$ \parencite{golub_matrix_2013} and define the whitened system
\begin{equation}
  \widetilde{\Mmat} = \mathbf{L}^{-1}\Mmat, \qquad
  \tilde{\dvec} = \mathbf{L}^{-1}\dvec ,
  \label{eq:whitening}
\end{equation}
so that $\Phi(\cvec) = \lVert \widetilde{\Mmat}\cvec - \tilde{\dvec}\rVert_2^2$: generalised
least squares becomes ordinary least squares on the whitened system, whose residuals have
identity covariance. When $\Cd$ is diagonal this reduces to dividing row $i$ of $\Mmat$ and
$d_i$ by $\sigma_i$.

\begin{remark}[Relation to the log-ratio objective]
\label{rem:logratio}
Suppose errors are proportional, $\sigma_i = f d_i$ for a constant relative error $f$ --- the
standard model for geochemical assay. The $i$th whitened residual is then
\begin{equation}
  \frac{(\Mmat\cvec - \dvec)_i}{f d_i}
  \;=\; \frac{1}{f}\left( \frac{(\Mmat\cvec)_i}{d_i} - 1 \right) ,
\end{equation}
which depends on the data only through the \textit{ratio} of predicted to observed
concentration. The weighted linear objective is therefore, like the log-ratio objective of the
parent study, a relative misfit; it is the unweighted objective that is not. This is precisely
the objection \textcite{barnes_using_2024} raise against Euclidean misfit --- that it
over-weights high concentrations --- and weighting by a proportional error model answers it.
The two remain distinct: \eqref{eq:gls} penalises the squared relative residual, whereas
$\max(x, 1/x)$ with $x = d_{\text{pred}}/d_{\text{obs}}$ behaves like $1 + |x - 1|$ near
$x = 1$. Both are functions of the relative residual alone, but one is quadratic and the other
piecewise linear in it, so they weight outliers differently and do not coincide.
\end{remark}

\section{Regularization}
\label{sec:regularization}

Section~\ref{sec:degenerate} showed that the unregularized problem interpolates the noise. The
remedy is Tikhonov regularization \parencite{tikhonov_solution_1963, hoerl_ridge_1970,
hansen_rank-deficient_1998}: add a penalty expressing a preference among models.

We penalise the \textbf{variance of the model about its own mean}. Let
\begin{equation}
  \Pmat \;=\; \Imat - \tfrac{1}{n}\onevec\onevec\transpose
  \label{eq:P}
\end{equation}
be the centering projector, so that $(\Pmat\cvec)_i = c_i - \bar{c}$. The full problem is
\begin{equation}
  \boxed{\;
  \begin{aligned}
    \underset{\cvec \in \Reals^n}{\text{minimize}} \quad
      & J(\cvec) = (\Mmat\cvec - \dvec)\transpose \Cd^{-1}(\Mmat\cvec - \dvec)
        + \lambda \lVert \Pmat\cvec \rVert_2^2 \\
    \text{subject to} \quad & \cvec \ge 0 .
  \end{aligned}
  \;}
  \label{eq:problem}
\end{equation}

\noindent Note that this penalises departure from the mean of the \textit{model}, not, as in
the parent study, from the mean of the \textit{observations}. That substitution was forced
there by the requirement that the objective be expressible in disciplined convex form
\parencite{boyd_convex_2004, diamond_cvxpy_2016}; the linear formulation is under no such
constraint and can use the more natural quantity.

\begin{lemma}[Properties of $\Pmat$]\label{lem:P}
$\Pmat$ is symmetric and idempotent, hence an orthogonal projector; it is positive
semi-definite with eigenvalue $0$ on $\operatorname{span}\{\onevec\}$ and eigenvalue $1$ on its
orthogonal complement; and
\begin{equation}
  \lVert \Pmat \cvec \rVert_2^2 = \cvec\transpose \Pmat\transpose \Pmat \cvec
  = \cvec\transpose \Pmat \cvec = \sum_{i=1}^n (c_i - \bar{c})^2 .
\end{equation}
\end{lemma}

\begin{proof}
Symmetry is clear. For idempotency, using $\onevec\transpose\onevec = n$,
\[
  \Pmat^2 = \Imat - \tfrac{2}{n}\onevec\onevec\transpose
          + \tfrac{1}{n^2}\onevec(\onevec\transpose\onevec)\onevec\transpose
          = \Imat - \tfrac{1}{n}\onevec\onevec\transpose = \Pmat .
\]
Then $\Pmat\transpose\Pmat = \Pmat^2 = \Pmat$, giving the quadratic form. $\Pmat\onevec =
\onevec - \onevec = \bm{0}$, and for $\mathbf{v} \perp \onevec$, $\Pmat\mathbf{v} =
\mathbf{v}$. Eigenvalues in $\{0,1\}$ give $\Pmat \succeq 0$.
\end{proof}

\noindent $\Pmat$ is thus singular, with a one-dimensional null space. This matters: the
penalty alone does not determine a unique solution, and the classical uniqueness condition for
general-form Tikhonov regularization \parencite{elden_algorithms_1977,
hansen_rank-deficient_1998} requires
$\mathcal{N}(\Mmat) \cap \mathcal{N}(\Pmat) = \{\bm{0}\}$. Here $\mathcal{N}(\Mmat) =
\{\bm{0}\}$ by Proposition~\ref{prop:invertible}, so the condition holds trivially --- but it holds
\textit{because of} the structure of $\Mmat$, not because of the penalty.

\begin{theorem}[Existence and uniqueness]\label{prop:unique}
For every $\lambda \ge 0$, $J$ is strictly convex on $\Reals^n$. Consequently
\eqref{eq:problem} has a unique global minimiser, both unconstrained and subject to
$\cvec \ge 0$.
\end{theorem}

\begin{proof}
$J$ is a sum of quadratic forms, with gradient and Hessian
\begin{align}
  \nabla J(\cvec) &= 2\Mmat\transpose\Cd^{-1}(\Mmat\cvec - \dvec) + 2\lambda\Pmat\cvec ,
  \label{eq:grad}\\
  \nabla^2 J &= 2\left(\Mmat\transpose\Cd^{-1}\Mmat + \lambda\Pmat\right) .
  \label{eq:hess}
\end{align}
For any $\mathbf{v} \ne \bm{0}$, $\mathbf{v}\transpose\Mmat\transpose\Cd^{-1}\Mmat\mathbf{v} =
(\Mmat\mathbf{v})\transpose\Cd^{-1}(\Mmat\mathbf{v}) > 0$, since $\Mmat\mathbf{v} \ne \bm{0}$
by Proposition~\ref{prop:invertible} and $\Cd^{-1} \succ 0$. Hence
$\Mmat\transpose\Cd^{-1}\Mmat \succ 0$. By Lemma~\ref{lem:P}, $\lambda\Pmat \succeq 0$ for
$\lambda \ge 0$. A positive definite matrix plus a positive semi-definite one is positive
definite \parencite{horn_matrix_2012}, so $\nabla^2 J \succ 0$ and $J$ is strictly convex. A
strictly convex function has at most one minimiser on any convex set, and coercivity
(guaranteed by $\nabla^2 J \succ 0$) ensures one exists. The non-negative orthant is convex and
closed, so the constrained problem inherits both.
\end{proof}

\begin{remark}
Uniqueness rests on $\Mmat$ being invertible, \textit{not} on the penalty. Had $\Mmat$ been
singular with a constant vector in its null space, $\lambda\Pmat$ could not have restored
uniqueness, because $\Pmat$ is blind to exactly that direction. Regularization here buys
conditioning and stability, not identifiability --- the problem was already identifiable.
\end{remark}

\begin{theorem}[The estimator]\label{prop:estimator}
Define
\begin{equation}
  \Amat_\lambda \;=\; \Mmat\transpose\Cd^{-1}\Mmat + \lambda\Pmat .
  \label{eq:A}
\end{equation}
If the unconstrained minimiser of $J$ is non-negative, it is also the solution of
\eqref{eq:problem}, and it is given by the \textbf{linear} map
\begin{equation}
  \boxed{\;
  \hat{\cvec} = \Rmat_\lambda \dvec ,
  \qquad
  \Rmat_\lambda = \Amat_\lambda^{-1}\Mmat\transpose\Cd^{-1} .
  \;}
  \label{eq:estimator}
\end{equation}
\end{theorem}

\begin{proof}
Setting \eqref{eq:grad} to zero gives the normal equations
$(\Mmat\transpose\Cd^{-1}\Mmat + \lambda\Pmat)\cvec = \Mmat\transpose\Cd^{-1}\dvec$, i.e.\
$\Amat_\lambda\cvec = \Mmat\transpose\Cd^{-1}\dvec$. $\Amat_\lambda \succ 0$ by the proof of
Proposition~\ref{prop:unique}, hence invertible, giving \eqref{eq:estimator}. For the first claim:
$J$ is strictly convex, so $\hat\cvec$ is its unique global minimiser over all of $\Reals^n$;
if it happens to be feasible then it is a fortiori the minimiser over the smaller feasible set.
\end{proof}

\noindent Two features of \eqref{eq:estimator} deserve emphasis. First, it is exactly linear in
\dvec, with \textbf{no constant offset}. This is a consequence of penalising
$\lVert\Pmat\cvec\rVert^2$, which is homogeneous in \cvec; a penalty
$\lVert\cvec - \cvec_0\rVert^2$ towards a fixed reference $\cvec_0 \ne \bm{0}$ would contribute
an affine term and complicate everything downstream. Second,
Proposition~\ref{prop:estimator} gives a practical recipe: solve the problem, and if the answer is
non-negative, report the analytical $\Rmat_\lambda\dvec$. The estimate and its covariance then
provably come from the same linear operator, which is what makes the error propagation of
Section~\ref{sec:propagation} meaningful. When the constraint \textit{is} active the estimator
is only piecewise linear, and that guarantee is lost.

\subsection{Limits}
\label{sec:limits}

\begin{theorem}[Small $\lambda$]
$\Rmat_0 = \Mmat^{-1}$, recovering Section~\ref{sec:inverse}.
\end{theorem}

\begin{proof}
Set $\lambda = 0$ in \eqref{eq:estimator}: $\Rmat_0 =
(\Mmat\transpose\Cd^{-1}\Mmat)^{-1}\Mmat\transpose\Cd^{-1} = \Mmat^{-1}$, as in
Proposition~\ref{prop:degenerate}.
\end{proof}

\begin{theorem}[Large $\lambda$]\label{prop:largelambda}
As $\lambda \to \infty$, $\hat{\cvec} \to t^\star\onevec$ with
\begin{equation}
  t^\star = \frac{\mathbf{m}\transpose\Cd^{-1}\dvec}{\mathbf{m}\transpose\Cd^{-1}\mathbf{m}},
  \qquad \mathbf{m} = \Mmat\onevec .
\end{equation}
If in addition the tracer is conservative and $\Cd = \sigma^2\Imat$, then
$t^\star = \bar{d}$, the arithmetic mean of the observations.
\end{theorem}

\begin{proof}
The penalty term $\lambda\lVert\Pmat\cvec\rVert^2$ diverges unless $\Pmat\cvec \to \bm 0$, i.e.\
unless \cvec approaches $\mathcal{N}(\Pmat) = \operatorname{span}\{\onevec\}$. Writing
$\cvec = t\onevec$ the penalty vanishes identically and the problem reduces to minimising
$(t\mathbf{m} - \dvec)\transpose\Cd^{-1}(t\mathbf{m} - \dvec)$ over the scalar $t$.
Differentiating and setting to zero gives $t^\star$ as stated. For a conservative tracer
Lemma~\ref{lem:stochastic} gives $\mathbf{m} = \Mmat\onevec = \onevec$, whence with
$\Cd = \sigma^2\Imat$ we get $t^\star = \onevec\transpose\dvec / \onevec\transpose\onevec =
\bar{d}$.
\end{proof}

\noindent So the regularization path runs from the exact, noise-interpolating inverse at
$\lambda = 0$ to a spatially uniform source field at $\lambda \to \infty$, whose level is the
(precision-weighted) mean observation. Both endpoints are interpretable, which is a useful
property when selecting $\lambda$ from the trade-off curve between misfit
$\lVert\Cd^{-1/2}(\Mmat\hat\cvec - \dvec)\rVert_2$ and roughness
$\lVert\Pmat\hat\cvec\rVert_2$ \parencite{hansen_rank-deficient_1998}.

\section{Error propagation}
\label{sec:propagation}

We now propagate the observational uncertainty $\Cd$ into the recovered sources. Because
$\hat\cvec$ is an exactly linear function of \dvec, this is not a linearisation: it is exact.

\begin{theorem}[Covariance of the sources]\label{prop:Cc}
If $\dvec$ has covariance $\Cd$ and no constraint is active, then
\begin{equation}
  \boxed{\;\Cc \;=\; \Rmat_\lambda\, \Cd\, \Rmat_\lambda\transpose . \;}
  \label{eq:Cc}
\end{equation}
\end{theorem}

\begin{proof}
For a deterministic matrix $\Rmat$ and random vector $\dvec$,
$\operatorname{Cov}(\Rmat\dvec) = \Rmat \operatorname{Cov}(\dvec) \Rmat\transpose$, directly
from the definition of covariance \parencite{menke_geophysical_2012}. Apply with
$\Rmat = \Rmat_\lambda$, which is deterministic given the network and the fluxes.
\end{proof}

\begin{theorem}[Covariance of the modelled observations]\label{prop:Cdhat}
The modelled observations $\hat\dvec = \Mmat\hat\cvec$ have
\begin{equation}
  \boxed{\;\Cdhat \;=\; \Mmat\, \Cc\, \Mmat\transpose
    \;=\; (\Mmat\Rmat_\lambda)\,\Cd\,(\Mmat\Rmat_\lambda)\transpose , \;}
  \label{eq:Cdhat}
\end{equation}
and at $\lambda = 0$, $\Cdhat = \Cd$ exactly.
\end{theorem}

\begin{proof}
The first equality is Proposition~\ref{prop:Cc} applied to the linear map $\hat\cvec \mapsto
\Mmat\hat\cvec$; the second follows by substitution. $\Mmat\Rmat_\lambda$ is the hat (influence)
matrix. At $\lambda = 0$, $\Rmat_0 = \Mmat^{-1}$ so $\Mmat\Rmat_0 = \Imat$ and $\Cdhat = \Cd$.
This is as it must be: by Proposition~\ref{prop:degenerate} the unregularized fit reproduces the
observations exactly, so the modelled observations inherit their uncertainty unchanged.
\end{proof}

\subsection{The unregularized case is efficient}

\begin{theorem}[Cram\'er--Rao]\label{prop:crb}
At $\lambda = 0$,
\begin{equation}
  \Cc = \Mmat^{-1}\Cd\Mmat^{-\mathsf{T}} = \left(\Mmat\transpose\Cd^{-1}\Mmat\right)^{-1} ,
\end{equation}
which is the inverse Fisher information of the model $\dvec \sim
\mathcal{N}(\Mmat\cvec, \Cd)$. The unregularized estimator therefore attains the Cram\'er--Rao
lower bound.
\end{theorem}

\begin{proof}
For the identity, note $(\Mmat^{-1}\Cd\Mmat^{-\mathsf{T}})^{-1} = \Mmat\transpose\Cd^{-1}\Mmat$
by inverting each factor and reversing the order. For the Fisher information, the
log-likelihood is $\ell(\cvec) = -\tfrac{1}{2}(\dvec - \Mmat\cvec)\transpose\Cd^{-1}(\dvec -
\Mmat\cvec) + \text{const}$, so $\nabla_{\cvec}\ell = \Mmat\transpose\Cd^{-1}(\dvec -
\Mmat\cvec)$ and $-\nabla^2_{\cvec}\ell = \Mmat\transpose\Cd^{-1}\Mmat$, which is
deterministic and hence equal to its expectation. The estimator is unbiased (see
Proposition~\ref{prop:bias} with $\lambda = 0$), so the Cram\'er--Rao bound applies and is met with
equality.
\end{proof}

\noindent No estimator, weighted or otherwise, can do better than $\Mmat^{-1}\dvec$ in the
unbiased class. Section~\ref{sec:degenerate} already told us weighting cannot help; this tells
us why nothing else can either. The only way to improve on \eqref{eq:varc} is to accept bias
--- which is exactly what regularization does.

\subsection{Bias, resolution, and the two covariances}

\begin{theorem}[Bias and resolution]\label{prop:bias}
$\Expect[\hat\cvec] = \Rmat_\lambda\Mmat\cvec_{\text{true}}$, so the estimator has bias
\begin{equation}
  \mathbf{b} = (\Rmat_\lambda\Mmat - \Imat)\cvec_{\text{true}}
             = -\lambda\,\Amat_\lambda^{-1}\Pmat\,\cvec_{\text{true}} ,
  \label{eq:bias}
\end{equation}
and \textbf{resolution matrix}
\begin{equation}
  \Rmat_\lambda\Mmat = \Imat - \lambda\Amat_\lambda^{-1}\Pmat .
  \label{eq:resolution}
\end{equation}
The mean-squared error is $\Expect[(\hat\cvec - \cvec_{\text{true}})(\cdot)\transpose]
= \Cc + \mathbf{b}\mathbf{b}\transpose$.
\end{theorem}

\begin{proof}
Substituting $\dvec = \Mmat\cvec_{\text{true}} + \bm{\varepsilon}$ with
$\Expect[\bm{\varepsilon}] = \bm 0$ into \eqref{eq:estimator} gives the first claim. Then
\[
  \Rmat_\lambda\Mmat - \Imat
  = \Amat_\lambda^{-1}\Mmat\transpose\Cd^{-1}\Mmat - \Imat
  = \Amat_\lambda^{-1}\!\left(\Mmat\transpose\Cd^{-1}\Mmat - \Amat_\lambda\right)
  = -\lambda\Amat_\lambda^{-1}\Pmat ,
\]
using \eqref{eq:A}. The MSE decomposition is the standard bias--variance identity
\parencite{hoerl_ridge_1970, aster_parameter_2018}.
\end{proof}

\noindent The bias vanishes when $\lambda = 0$, and also whenever $\Pmat\cvec_{\text{true}} =
\bm 0$, i.e.\ when the true source field really is uniform --- the penalty is then telling the
truth and costs nothing. Since $\cvec_{\text{true}}$ is unknown, $\mathbf{b}$ cannot be
evaluated; the resolution matrix \eqref{eq:resolution} is the standard computable proxy
\parencite{menke_geophysical_2012}, equal to $\Imat$ at $\lambda = 0$ and degrading as
$\lambda$ grows. Its trace, $\operatorname{tr}(\Rmat_\lambda\Mmat)$, is the effective number of
degrees of freedom the data constrain.

This bias is precisely why \eqref{eq:Cc} must be interpreted with care. It answers the
question ``if I repeated the measurements, how would my estimate scatter?'' --- and a heavily
damped estimator scatters very little while being systematically wrong. A second covariance
answers a different and often more useful question.

\begin{theorem}[Posterior covariance]\label{prop:posterior}
Read the penalty as a Gaussian prior with precision $\lambda\Pmat$ and the likelihood as
$\dvec \mid \cvec \sim \mathcal{N}(\Mmat\cvec, \Cd)$. Then the posterior covariance is
$\Amat_\lambda^{-1}$, and
\begin{equation}
  \boxed{\;
  \Rmat_\lambda\Cd\Rmat_\lambda\transpose
  \;=\; \Amat_\lambda^{-1} \;-\; \lambda\,\Amat_\lambda^{-1}\Pmat\Amat_\lambda^{-1} .
  \;}
  \label{eq:identity}
\end{equation}
Consequently $\Cc \preceq \Amat_\lambda^{-1}$, with equality if and only if $\lambda = 0$.
\end{theorem}

\begin{proof}
For the linear-Gaussian model the posterior precision is the sum of the likelihood precision
$\Mmat\transpose\Cd^{-1}\Mmat$ and the prior precision $\lambda\Pmat$, giving $\Amat_\lambda$
\parencite{tarantola_inverse_2005}. For the identity, substitute \eqref{eq:estimator} and use
$\Cd^{-1}\Cd\Cd^{-1} = \Cd^{-1}$:
\begin{align*}
  \Rmat_\lambda\Cd\Rmat_\lambda\transpose
  &= \Amat_\lambda^{-1}\Mmat\transpose\Cd^{-1}\,\Cd\,\Cd^{-1}\Mmat\Amat_\lambda^{-1} \\
  &= \Amat_\lambda^{-1}\left(\Mmat\transpose\Cd^{-1}\Mmat\right)\Amat_\lambda^{-1} \\
  &= \Amat_\lambda^{-1}\left(\Amat_\lambda - \lambda\Pmat\right)\Amat_\lambda^{-1}
   = \Amat_\lambda^{-1} - \lambda\Amat_\lambda^{-1}\Pmat\Amat_\lambda^{-1},
\end{align*}
using the symmetry of $\Amat_\lambda$. For the ordering, note that for any $\mathbf{x}$,
$\mathbf{x}\transpose\Amat_\lambda^{-1}\Pmat\Amat_\lambda^{-1}\mathbf{x} =
(\Amat_\lambda^{-1}\mathbf{x})\transpose\Pmat(\Amat_\lambda^{-1}\mathbf{x}) \ge 0$ by
Lemma~\ref{lem:P}, so the subtracted term is positive semi-definite. It is zero for all
$\mathbf{x}$ only if $\lambda = 0$, since $\Amat_\lambda^{-1}$ is non-singular and $\Pmat \ne
\bm 0$.
\end{proof}

\noindent The propagated covariance is therefore always the smaller of the two, and quoting it
alone at large $\lambda$ understates the true error. The posterior covariance
$\Amat_\lambda^{-1}$ accounts for what the damping costs as well as what it stabilises, and is
the more honest error bar to report whenever $\lambda > 0$. At $\lambda = 0$ the two coincide
and both equal the Cram\'er--Rao bound of Proposition~\ref{prop:crb}.

\section{Summary}

\begin{table}[h]
\centering
\small
\begin{tabular}{@{}ll@{}}
\toprule
Quantity & Expression \\
\midrule
Mixing matrix & $M_{ij} = \Theta_{ij}\alpha_{j\to i} q_j / Q_i$ \\
Exact inverse & $c_i = \left(Q_i d_i - \sum_{p \in A(i)}\alpha_{p\to i}Q_p d_p\right)/q_i$ \\
Amplification & $\kappa_i = Q_i/q_i = 1/M_{ii}$ \\
Penalty operator & $\Pmat = \Imat - n^{-1}\onevec\onevec\transpose$ \\
Normal matrix & $\Amat_\lambda = \Mmat\transpose\Cd^{-1}\Mmat + \lambda\Pmat$ \\
Estimator & $\Rmat_\lambda = \Amat_\lambda^{-1}\Mmat\transpose\Cd^{-1}$ \\
Source covariance & $\Cc = \Rmat_\lambda\Cd\Rmat_\lambda\transpose$ \\
Prediction covariance & $\Cdhat = \Mmat\Cc\Mmat\transpose$ \\
Posterior covariance & $\Amat_\lambda^{-1}$ \\
Resolution & $\Rmat_\lambda\Mmat = \Imat - \lambda\Amat_\lambda^{-1}\Pmat$ \\
Bias & $\mathbf{b} = -\lambda\Amat_\lambda^{-1}\Pmat\cvec_{\text{true}}$ \\
\bottomrule
\end{tabular}
\caption{\textbf{Principal results.} At $\lambda = 0$: $\Rmat_0 = \Mmat^{-1}$, the resolution
matrix is $\Imat$, the bias vanishes, the two covariances coincide, and $\Cc$ attains the
Cram\'er--Rao bound.}
\end{table}

\noindent The practical picture is this. The mixing matrix is invertible, so the unmixing
problem always has an exact solution --- but that solution differences neighbouring
observations, amplifying their errors by $\kappa_i = Q_i/q_i$, and interpolates the noise.
Regularization trades this variance for bias, and because the penalty is a homogeneous
quadratic the resulting estimator stays exactly linear in the data. That linearity is what
delivers closed-form covariances, a resolution matrix, and effective degrees of freedom, none
of which are available from the log-ratio formulation without Monte Carlo resampling. The price
is the assumption of absolute rather than relative misfit --- an assumption that
Remark~\ref{rem:logratio} shows is substantially repaired by weighting with a proportional
error model, and which is in any case mild for tracers of low log-variance such as isotopic
ratios \parencite{blondes_practical_2016}. For tracers whose information really is relative
rather than absolute --- that is, genuinely compositional data
\parencite{aitchison_statistical_1986} --- the log-ratio formulation remains the better choice.

Two caveats bear repeating. Regularization is not optional: by
Proposition~\ref{prop:degenerate} the unregularized problem reproduces the data exactly, noise
included. And the analytical machinery of Section~\ref{sec:propagation} is valid only while
the constraint $\cvec \ge 0$ is inactive; when sites clamp, the estimator ceases to be linear
and the reported covariances describe the unconstrained estimator instead.

\printbibliography

\end{document}
